% sections/intro.tex -- Introduction (final, W19). Every number is a macro from
% numbers.tex; the defence paragraph waits for W17 (E7/E8).
\section{Introduction}\label{sec:intro}

Sender-constrained access tokens are the standard defence against token
theft in high-assurance APIs. FAPI~2.0 requires them, and RFC~9449 (DPoP) and
RFC~8705 provide them~\cite{rfc9449,rfc8705,fapisp2025}. Formal analyses prove the
resulting protocols secure, \emph{provided} the resource server performs every
specified check~\cite{hosseyni2024csf,hosseyni2024tops}.

Those analyses model the resource server as a single party. Real deployments
split it. A gateway plugin or an external authorization service checks the
proof against \emph{its view} of the request; a framework router,
method-override middleware and a handler decide what \emph{actually runs}.
DPoP binds only the method and the target URI without its
query~\cite{rfc9449}. Whenever the two components' views of a request
diverge, or the executor decides on something the proof does not cover, the
sender constraint degrades without any check failing.

We make this precise as \emph{binding soundness}: a proof accepted for the
verifier's view must determine exactly one executed operation
(Definition~\ref{def:soundness}). Its \emph{binding resolution} counts how
many operations one accepted proof can select. Resolution exceeds one in
three ways: \GA, under-binding by design, where the query, the body or a method
override selects the operation; \GB, interpretation differentials between
verifier and router; and \GC, spoofable propagation of check results between
hops. A related property, \emph{authorization soundness}, asks the same
question of a policy decision made on the verifier's view
(Definition~\ref{def:authz}). It matters for a malicious but legitimate
client, who can sign any request it likes.

We test both properties with a differential harness whose oracle is the
definition itself: the verifier accepts the request, and the executor runs a
different operation, judged against ground truth measured on the bare
executor. The harness combines \NumGateways{} gateways, \NumFrameworks{}
frameworks implementing one bench API, a single reference verifier and
\NumPrimitives{} mutation primitives sent as raw bytes. We ran
\NumTopologiesAll{} topologies and \NumCasesAll{} requests, and re-ran every
finding three times on fresh deployments. A prevalence study of
\NumSpecs{} API descriptions and the override code of \NumTargetFrameworks{}
frameworks estimates how often \GA occurs outside the harness.

\paragraph{Findings.} Two interpretation differentials carry the result.
First, under a malicious client with a path policy at the gateway, we
found \NumEFourViolations{} cross-tenant or admin executions in
\NumEFourClasses{} classes. Two of them, Kestrel's and Kong's handling of
\texttt{/../}, defeat a strict default-deny allow-list, Kong's for every
backend. Second, under a leaked proof, NGINX and Kong forward a \texttt{\#}
that three backend servers keep in the path. With a nested route, a leaked
\texttt{deleteAccount} proof then runs \texttt{deleteStatement} under default
configurations. Both rest on component behaviours that are individually
documented or already known; what is new is that they break the binding
between a correctly verified proof and the operation it authorises.
Overall, \NumVulnTopologies{} of \NumTopologies{} topologies have binding
resolution above one under a leaked proof (up to \MaxResolution{}), with
\NumHits{} confirmed hits. Most are \GA: query and body dispatch change the
operation behind every gateway and framework, and with method override
enabled a POST proof executes PUT or DELETE in all \NumFrameworks{}
frameworks, by default in Rails and Laravel. We also observed an
encoded-slash split behind APISIX whose mechanism we have not yet confirmed,
and a low-impact dot-segment effect on catch-all routes. We found no \GC
binding hit. Verifying the same proof again at the executor removes none of
the \GA hits and not the \texttt{\#} differential.

Two defences that check the executor's routed view (D-A) or a view MAC it
recomputes (D-B) remove the override and all \GB witnesses, leaving the
\NumDefResidualDA{} query and body dispatch witnesses of \NumDefWitnesses{}, which need
query and body binding (\Abl{1}). D-B costs \DefenceOverheadPfifty{} at the
median, but rejects every legitimate request whose path a gateway rewrote.

\paragraph{Contributions.}
\begin{itemize}[leftmargin=*,itemsep=1pt]
\item \textbf{Property.} Binding soundness and binding resolution for
sender-constrained requests in layered deployments, separated from
authorization soundness, with three failure classes and two adversaries
(\S\ref{sec:model}).
\item \textbf{Measurement of \GA.} How often API descriptions put several
state-changing operations behind one (method, path) key, and where method
override is on by default (\S\ref{sec:prevalence}).
\item \textbf{A binding oracle and its findings.} An open harness that
measures proof acceptance against the executed operation across gateways,
placements and frameworks, with the root causes it found
(\S\ref{sec:harness}).
\item \textbf{Recommendations.} An executor-view contract for the standard,
and for gateways, authorizing on the representation they forward
(\S\ref{sec:defence}, \S\ref{sec:discussion}).
\end{itemize}

All experiments ran on local containers; no third-party system was tested
(\S\ref{sec:discussion:ethics}).
